\subsection{The Holy Night}\label{sec:holy-night}

Advent ends on the Vigil at None, the ninth hour, the short hour of the Office
prayed in the middle of the afternoon. With First Vespers that evening the
Church is already keeping Christmas.

If you keep the Vigil as a day of fasting, let it be the last day of the fast.
Finish what must be finished early, so that the evening is free. If there is a
crib in the house, set it up during the day; you may leave the manger empty
until night, and let the youngest member of the household lay the Child in it
when the time comes. Read the Martyrology's proclamation of the birth once more
beside it, and kneel at the words of the birth.

At First Vespers the Church sings a new hymn, \latin{Jesu Redemptor omnium},
whose last stanza begins, ``O Jesu! born of Virgin bright, Immortal glory be
to Thee''. The O antiphons are finished, and the \latin{Rorate} and the
\latin{Excita} with them. What the Church asked for has come.

\subsubsection*{At the Manger}

The station of the First Sunday was St Mary Major, and on the Vigil the Church
went back there. The first Mass of Christmas, celebrated at night, has its
station at St Mary Major again, with two words added: \latin{ad Praesepe}, at
the Manger. On the holy night the season goes, as it were, one step further in,
to the Manger itself. The Mass at dawn is at St Anastasia, and the Mass of the
day at St Mary Major once more.

St Gregory the Great, preaching in the basilica of St Mary on Christmas Day, began by
excusing himself for being brief: ``By God's mercy we are to say three Masses
to-day, so that there is not much time left for preaching; but at the same time
the occasion of the Lord's Birth-day itself obliges me to speak a few words.''
In those few words he spoke of the census, of Bethlehem the House of Bread, and
of the birth away from home.

Go to the Mass at night if you can. Before it, or at home if you cannot go,
read St Luke's account of the birth aloud once more. Then come to the crib. The
ox and the ass are there, that know their master's crib; the Child is there,
wrapped in swaddling bands and laid in the place where beasts are fed. At
Lauds of Christmas the Church sings:

\begin{verse}
A manger scantly strewn with hay \\
Becomes th' Eternal's bed; \\
And He, who feeds each smallest bird, \\
Himself with milk is fed.
\end{verse}

\subsubsection*{Let us rejoice}

At Matins of Christmas the Church reads St Leo the Great. He begins with the
angel's words, ``Unto us is born this day a Saviour'', and with ``Let us
rejoice''; then he forbids sadness:

\begin{quote}
It would be unlawful to be sad to-day, for to-day is Life's Birthday: the
Birthday of that Life, Which, for us dying creatures, taketh away the sting of
death, and bringeth the bright promise of the eternal gladness hereafter. It
would be unlawful for any man to refuse to partake in our rejoicing. All men
have an equal share in the great cause of our joy, for, since our Lord, Who is
the destroyer of sin and of death, findeth that all are bound under the
condemnation, He is come to make all free. Rejoice, O thou that art holy, thou
drawest nearer to thy crown! Rejoice, O thou that art sinful, thy Saviour
offereth thee pardon! Rejoice also, O thou Gentile, God calleth thee to life!\par
{\smallskip\noindent\small ---St Leo the Great, Sermon 21 on the Nativity\par}
\end{quote}

The \latin{Gloria in excelsis}, which Durandus said the Church holds back in
Advent so that it may be sung the more eagerly on this night, is the angels' own
song, and tonight it is sung again: ``Glory to God in the highest: and on earth
peace to men of good will.''

St Gregory Nazianzen began his sermon on the Birthday of Christ with a cry
that the whole of Advent has been waiting to make:

\begin{quote}
Christ is born, glorify ye Him. Christ from heaven, go ye out to meet Him.
Christ on earth; be ye exalted. Sing unto the Lord all the whole earth; and
that I may join both in one word, Let the heavens rejoice, and let the earth be
glad, for Him Who is of heaven and then of earth. Christ in the flesh, rejoice
with trembling and with joy; with trembling because of your sins, with joy
because of your hope.\par
{\smallskip\noindent\small ---St Gregory Nazianzen, Oration 38, on the Birthday of Christ\par}
\end{quote}
